
% ══════════════════════════════════════════════════════════════════════════════
%  APPENDIX A: DECLARATION OF AI ASSISTANCE
% ══════════════════════════════════════════════════════════════════════════════

\chapter{Declaration of AI Assistance}

\section{Use of Artificial Intelligence in This Document}

This course guide was developed with assistance from artificial intelligence (AI)
in the following capacities:

\begin{itemize}
  \item \textbf{Template Design:} AI provided guidance on structuring the document layout,
        choosing appropriate color schemes.

  \item \textbf{LaTeX Code Generation:} AI assisted in writing and refactoring LaTeX code,
        including package configurations, box styling (tcolorbox), TikZ decorations,
        and custom command definitions.

  \item \textbf{LaTeX Structure Refactoring:} AI helped reorganize and optimize the document
        structure for maintainability, readability, and professional presentation.

  \item \textbf{Mathematical Formula Formatting:} AI provided assistance in formatting and
        presenting mathematical equations in clear, readable notation aligned with
        conventions.
\end{itemize}

\bigskip\noindent
\textbf{Author Verification:} All content, including scientific accuracy, theoretical derivations,
example problems, and pedagogical choices, has been carefully reviewed and verified by the author.
The AI assistance was limited to technical formatting and structural optimization, not to
the generation or validation of course content itself.

\bigskip\noindent
\textit{Compiled by: \NoteAuthor \quad | \quad Date: \CompileDate}

\clearpage

% ══════════════════════════════════════════════════════════════════════════════
%  APPENDIX B: MATHEMATICAL TOOLS & DERIVATIONS
% ══════════════════════════════════════════════════════════════════════════════

\chapter{Mathematical Tools \& Derivations}

\section{Useful Identities}

\subsection*{Trigonometric}
$\sin^2\theta + \cos^2\theta = 1$,\quad
$\tan\theta = \sin\theta/\cos\theta$,\quad
$\sin(A\pm B) = \sin A\cos B \pm \cos A\sin B$.

\subsection*{Hyperbolic}
$\cosh^2 x - \sinh^2 x = 1$,\quad
$\sinh(x\pm y) = \sinh x\cosh y \pm \cosh x\sinh y$.

\subsection*{Common Integrals}
\begin{align*}
  \int x^n\,dx &= \tfrac{x^{n+1}}{n+1} + C \quad (n\neq-1), &
  \int e^{ax}\,dx &= \tfrac{1}{a}e^{ax}+C,\\
  \int \frac{dx}{x} &= \ln|x|+C, &
  \int \sin x\,dx &= -\cos x + C.
\end{align*}

\section{Differential Geometry Primer}

The \textbf{metric tensor} $g_{\mu\nu}$ encodes the geometry of spacetime:
$ds^2 = g_{\mu\nu}\,dx^\mu\,dx^\nu$.
The \textbf{Christoffel symbols} (not tensors) are defined as
\begin{equation}
  \Gamma^\lambda_{\mu\nu} = \tfrac{1}{2}g^{\lambda\sigma}
    (\partial_\mu g_{\nu\sigma} + \partial_\nu g_{\mu\sigma} - \partial_\sigma g_{\mu\nu}).
\end{equation}
The \textbf{Ricci tensor} is the contraction $R_{\mu\nu} = R^\lambda{}_{\mu\lambda\nu}$
of the Riemann curvature tensor, and the Ricci scalar is $R = g^{\mu\nu}R_{\mu\nu}$.
These enter the Einstein field equations as
$G_{\mu\nu} \equiv R_{\mu\nu} - \tfrac{1}{2}g_{\mu\nu}R = (8\pi G/c^4)T_{\mu\nu}$.

\clearpage

% ══════════════════════════════════════════════════════════════════════════════
%  APPENDIX C: CONSTANTS & CONVERSION FACTORS
% ══════════════════════════════════════════════════════════════════════════════

\chapter{Constants \& Conversion Factors}

\section{Physical Constants}

\begin{center}
\begin{tabular}{lll}
\toprule
\textbf{Quantity} & \textbf{Symbol} & \textbf{Value (SI)} \\
\midrule
Speed of light         & $c$       & $2.998 \times 10^8 \,\mathrm{m\,s^{-1}}$           \\
Gravitational constant & $G$       & $6.674 \times 10^{-11} \,\mathrm{m^3\,kg^{-1}\,s^{-2}}$ \\
Hubble constant today  & $H_0$     & $67.4 \pm 0.5 \,\mathrm{km\,s^{-1}\,Mpc^{-1}}$     \\
\bottomrule
\end{tabular}
\end{center}

\section{Conversion Factors}

\begin{center}
\begin{tabular}{ll}
\toprule
\textbf{Conversion} & \textbf{Value} \\
\midrule
1 parsec (pc)       & $3.086\times10^{16}\,\mathrm{m}$ \\
1 kiloparsec (kpc)  & $3.086\times10^{19}\,\mathrm{m}$ \\
1 megaparsec (Mpc)  & $3.086\times10^{22}\,\mathrm{m}$ \\
1 gigaparsec (Gpc)  & $3.086\times10^{25}\,\mathrm{m}$ \\
1 light-year (ly)   & $9.461\times10^{15}\,\mathrm{m}$ \\
1 gigalight-year (Gly) & $9.461\times10^{24}\,\mathrm{m}$ \\
1 year (yr)         & $3.156\times10^{7}\,\mathrm{s}$ \\
1~Gyr               & $3.156\times10^{16}\,\mathrm{s}$ \\
\bottomrule
\end{tabular}
\end{center}

\clearpage
